\documentclass{article}
\usepackage{amsmath, amssymb, amsthm}
\usepackage[a4paper,
            bindingoffset=0.2in,
            left=1in,
            right=1in,
            top=1in,
            bottom=1in,
            footskip=.25in]{geometry}
\newtheorem{theorem}{Teorema}[section]
\newtheorem{corollary}{Corolario}[theorem]
\newtheorem{lemma}[theorem]{Lema}
\renewcommand\qedsymbol{$\blacksquare$}
\renewenvironment{proof}{{\bfseries Dem:}}

\title{Ejercicios}
\author{Jose Hernandez}
\date{\today}

\begin{document}
\maketitle

\section{Parcial 2}

Sean \( G \) y \( G' \) grupos, y sea \( \psi  \colon G \to G' \). 

Definiciones importantes:

\begin{enumerate}
    \item Si \( \psi \) es tal que
        \[
            \psi (ab) = \psi (a) \psi (b)
        \] 
        entonces \( \psi  \) es un homomorfismo.
    \item Si \( \psi  \) es homomorfismo y ademas es inyectivo, entonces se llama monomorfismo.
    \item Si \( \psi  \) es monomorfismo y ademas es sobreyectivo, entonces se llama isomorfismo.
    \item Sea \( N \) subgrupo de \( G \). Se dice que \( N \) es un subgrupo normal de \( G \), si
        \( a^{-1}Na \subseteq N \), para todo \( a \in G \). y lo denotamos por \( N \triangleleft G \)
    \item Un grupo tal que no es abeliano, pero todos sus subgrupos son normales es llamado Grupo Hamiltoniano
\end{enumerate}


\begin{theorem}
    Si \( \psi  \) es un homomorfismo, entonces

    a) \( Ker(\psi ) \) es un subgrupo de \( G \)

    b) Sea \( a \in G \), \( a^{-1}(Ker(\psi ))a \subseteq Ker(\psi ) \)
\end{theorem}

\begin{proof}  

    a) Sean \( a,b \in  Ker(\psi) \), entonces, sabemos que \( \psi (a) = e' \) y \( \psi (b) = e' \).
    Notemos que 
    \begin{align*}
        \psi (ab^-1) &= \psi (a)(\psi (b))^{-1} \\  
                     &= e'(e')^{-1} \\
                     &= e'
    \end{align*}

    Asi, \( Ker(\psi ) \) es un subgrupo de \( G \).

    b) Sea \( x \in  a^{-1}(Ker(\psi ))a \), entonces, \( x = a^{-1}ba \) para algun \( b \in Ker(\psi)\),
    queremos ver que \( \psi (x) = e' \).

    \begin{align*}
        \psi (x) &= \psi (a^{-1}ba) \\
                 &= (\psi (a))^{-1}\psi (b)\psi (a) \\
                 &= (\psi (a))^{-1} \cdot e' \cdot \psi (a) \tag{por \( b \in Ker(\psi ) \)} \\
                 &= e'
    \end{align*}

    Asi, se obtiene el resultado buscado.
    \( \hfill \blacksquare \)

\end{proof}

\begin{corollary}
    Sea \( \psi  \) un homomorfismo. \( \psi  \) será un monomorfismo si, y solo si, \( Ker(\psi ) = \{e\}   \)
\end{corollary}

\begin{proof}
    \vspace{3mm}
    
    \( (\implies) \)

    Supongamos que \( \psi  \) es un monomorfismo.

    Veamos que \( \{e\} \supseteq Ker(\psi)  \).

    Sea \( a \in Ker(\psi) \), entonces 
    \begin{align*}
        \psi (a) &= e' = \psi (e) \\
        \psi (a) &= \psi (e)
    \end{align*}
    \qquad Luego, como \( \psi  \) es monomorfismo, se tiene que \( a=e \), es decir, \( Ker(\psi ) \subseteq \{e\}   \).
    \vspace{3mm}

    \( \left( \impliedby \right)  \)

    Supongamos que \( Ker(\psi ) = \{e\} \) y \( \psi  \) homomorfismo.

    Sean \( a,b \in G \) tales que \( \psi (a) = \psi (b) \), luego
    
    \begin{align*}
        \psi (aa^{-1}) &= \psi (a)(\psi (a))^{-1} \\
                       &= \psi (b)(\psi (a))^{-1} \\
                       &= \psi (ba^{-1}) \\
                       &= e' \\
                       &\iff \\
        ba^{-1} &= e  \tag{por hipótesis} \\
        b &= a
    \end{align*}
    Así, \( \psi  \) es monomorfismo.  \( \hfill \blacksquare \)

\end{proof}

\subsection*{Problemas pag. 73 Hernstein}
\subsection*{1} 
\begin{itemize}
    \item[a)] Si es homomorfismo y \( Ker(\varphi) = \{\text{todos los múltiplos de \( n \)}\}  \).

        No es 1-1, ya que \( \varphi (0) = [0] \) y \( \varphi (n) = [0]\), pero \( 0 \neq n \).

        Si es sobreyectiva, ya que todas las clases de equivalencia en \( \mathbb{Z}_{n} \) tienen al menos
        un representante en \( \mathbb{Z} \).

    \item[b)] No es un homomorfismo. \( \varphi (a b) = \varphi (b) \varphi (a) \).

    \item[c)] Si es un homomorfismo, ya que es el caso (b) pero conmutativo.

        Si es 1-1, ya que \( a^{-1} = b^{-1} \iff a = b \).

        Es sobreyectivo porque todo elemento en un grupo tiene su inverso en el mismo grupo.

    \item[d)] Si es homomorfismo

        \begin{equation*}
            \varphi (p \cdot q) = 
            \begin{cases}
                1  \text{si } pq > 0\\
                -1
            \end{cases}
            \text{pero esto }
        \end{equation*}
\end{itemize}

\end{document}
